\documentclass[10pt,a4paper]{amsart}
\usepackage{amsmath,amssymb,amsthm,graphicx}
% Retained preamble for the Stacks Project chapter morphisms.tex (Morphisms of Schemes).
% The chapter file pinned-original-source.tex carries only \input{preamble} on its first line, so no
% preamble text lives inside the chapter itself.  The Stacks Project uses one single shared preamble for
% every chapter; the snapshot stages it verbatim in sources/stacks-chow/preamble.tex.  The block below is
% copied mechanically, line for line, from that file's lines 14-18 (the Xy-pic setup) and lines 154-217
% (the native theorem declarations together with the mathematical macros).  Nothing was added, removed or
% redefined.  The wrapper class and the amsmath/amssymb/amsthm/graphicx base packages are supplied by the
% assembler, so the class/\externaldocument part of the shared preamble is not repeated here.
% For commutative diagrams we use Xy-pic
\usepackage[all]{xy}

% We use 2cell for 2-commutative diagrams.
\xyoption{2cell}
% Theorem environments.
%
\theoremstyle{plain}
\newtheorem{theorem}[subsection]{Theorem}
\newtheorem{proposition}[subsection]{Proposition}
\newtheorem{lemma}[subsection]{Lemma}

\theoremstyle{definition}
\newtheorem{definition}[subsection]{Definition}
\newtheorem{example}[subsection]{Example}
\newtheorem{exercise}[subsection]{Exercise}
\newtheorem{situation}[subsection]{Situation}

\theoremstyle{remark}
\newtheorem{remark}[subsection]{Remark}
\newtheorem{remarks}[subsection]{Remarks}

\numberwithin{equation}{subsection}

% Macros
%
\def\lim{\mathop{\mathrm{lim}}\nolimits}
\def\colim{\mathop{\mathrm{colim}}\nolimits}
\def\Spec{\mathop{\mathrm{Spec}}}
\def\Hom{\mathop{\mathrm{Hom}}\nolimits}
\def\Ext{\mathop{\mathrm{Ext}}\nolimits}
\def\SheafHom{\mathop{\mathcal{H}\!\mathit{om}}\nolimits}
\def\SheafExt{\mathop{\mathcal{E}\!\mathit{xt}}\nolimits}
\def\Sch{\mathit{Sch}}
\def\Mor{\mathop{\mathrm{Mor}}\nolimits}
\def\Ob{\mathop{\mathrm{Ob}}\nolimits}
\def\Sh{\mathop{\mathit{Sh}}\nolimits}
\def\NL{\mathop{N\!L}\nolimits}
\def\CH{\mathop{\mathrm{CH}}\nolimits}
\def\proetale{{pro\text{-}\acute{e}tale}}
\def\etale{{\acute{e}tale}}
\def\QCoh{\mathit{QCoh}}
\def\Ker{\mathop{\mathrm{Ker}}}
\def\Im{\mathop{\mathrm{Im}}}
\def\Coker{\mathop{\mathrm{Coker}}}
\def\Coim{\mathop{\mathrm{Coim}}}

% Boxtimes
%
\DeclareMathSymbol{\boxtimes}{\mathbin}{AMSa}{"02}

%
% Macros for moduli stacks/spaces
%
\def\QCohstack{\mathcal{QC}\!\mathit{oh}}
\def\Cohstack{\mathcal{C}\!\mathit{oh}}
\def\Spacesstack{\mathcal{S}\!\mathit{paces}}
\def\Quotfunctor{\mathrm{Quot}}
\def\Hilbfunctor{\mathrm{Hilb}}
\def\Curvesstack{\mathcal{C}\!\mathit{urves}}
\def\Polarizedstack{\mathcal{P}\!\mathit{olarized}}
\def\Complexesstack{\mathcal{C}\!\mathit{omplexes}}
% \Pic is the operator that assigns to X its picard group, usage \Pic(X)
% \Picardstack_{X/B} denotes the Picard stack of X over B
% \Picardfunctor_{X/B} denotes the Picard functor of X over B
\def\Pic{\mathop{\mathrm{Pic}}\nolimits}
\def\Picardstack{\mathcal{P}\!\mathit{ic}}
\def\Picardfunctor{\mathrm{Pic}}
\def\Deformationcategory{\mathcal{D}\!\mathit{ef}}

% --- appended 2026-10-12 (Chief): the project own definitions of the reference/slogan/history
% display environments, copied mechanically from the shared Stacks preamble lines 8-12
% (\usepackage{verbatim} plus the three \newenvironment lines, which the chapter preamble keeps
% verbatim from the project preamble). Nothing added, removed or redefined: these environments
% carry no mathematical content and the project itself typesets them as comments. Without these
% three lines a counted statement whose own text uses them fails with Environment slogan undefined.
\usepackage{verbatim}
\newenvironment{reference}{\comment}{\endcomment}
\newenvironment{slogan}{\comment}{\endcomment}
\newenvironment{history}{\comment}{\endcomment}
\usepackage{cleveref}
\crefname{originallemma}{original lemma}{original lemmas}
\Crefname{originallemma}{Original lemma}{Original lemmas}
\makeatletter
\newlabel{lemma-finite-proper}{{01WN}{source}{}{sourcepointer-1}{}}
\newlabel{lemma-finite-proper@cref}{{[originallemma][0][]01WN}{[1][0][]source}{}{}{}}
\makeatother
\crefname{originallemma}{original lemma}{original lemmas}
\Crefname{originallemma}{Original lemma}{Original lemmas}
\makeatletter
\newlabel{lemma-generically-finite}{{02NW}{source}{}{sourcepointer-2}{}}
\newlabel{lemma-generically-finite@cref}{{[originallemma][0][]02NW}{[1][0][]source}{}{}{}}
\makeatother
\crefname{originallemma}{original lemma}{original lemmas}
\Crefname{originallemma}{Original lemma}{Original lemmas}
\makeatletter
\newlabel{lemma-image-proper-scheme-closed}{{01W6}{source}{}{sourcepointer-3}{}}
\newlabel{lemma-image-proper-scheme-closed@cref}{{[originallemma][0][]01W6}{[1][0][]source}{}{}{}}
\makeatother
\crefname{originallemma}{original lemma}{original lemmas}
\Crefname{originallemma}{Original lemma}{Original lemmas}
\makeatletter
\newlabel{lemma-quasi-compact-dominant}{{01RL}{source}{}{sourcepointer-4}{}}
\newlabel{lemma-quasi-compact-dominant@cref}{{[originallemma][0][]01RL}{[1][0][]source}{}{}{}}
\makeatother
\crefname{originallemma}{original lemma}{original lemmas}
\Crefname{originallemma}{Original lemma}{Original lemmas}
\makeatletter
\newlabel{properties-lemma-quasi-compact-dense-open-separated}{{03J1}{source}{}{sourcepointer-5}{}}
\newlabel{properties-lemma-quasi-compact-dense-open-separated@cref}{{[originallemma][0][]03J1}{[1][0][]source}{}{}{}}
\makeatother
\title{A generically finite morphism over a finite generic fibre is finite over a dense open of the target}
\author{The Stacks Project Authors}
\date{stacks-project:morphisms.tex:local-snapshot-20261012; GNU FDL 1.2 or later}
\begin{document}
\maketitle
\noindent\textit{Editorial note.} All delivered mathematical text is the author's own, copied from the pinned Stacks Project chapter file pinned-original-source.tex; no statement or proof step was rewritten, reordered or removed.\par\medskip
\noindent\textit{Editorial note.} Following the multi-claim convention for this source, source\_url is the source-level URL https://stacks.math.columbia.edu/ and the source label lemma-finite-over-dense-open is carried in claim\_key and here rather than in a per-claim tag URL.\par\medskip
\noindent\textit{Editorial note.} The preamble is the chapter's shared Stacks Project preamble, retained as sources/stacks-morphisms/retained-preamble.tex: a line-for-line excerpt of the shared Stacks preamble staged with this snapshot. No macro was invented or redefined, and this card adds no preamble text of its own.\par\medskip
\noindent\textit{Editorial note.} Statement span 13328-13339 and proof span 13341-13367 come from the halfA pool card operations/pool-cards/stacks-morphisms-halfA.json and were re-read in the pinned chapter file before this card was written.\par\medskip
\noindent\textit{Editorial note.} The proof refers to the two alternatives in the hypothesis and treats them in the separated and quasi-compact cases; the author's wording is preserved.\par\medskip
\noindent\textit{Editorial note.} 5 cross-references occur inside the counted spans; they are listed in external\_references as navigation pointers only.\par\medskip
\noindent\textit{Allowed standard prerequisites (editorial).}\begin{itemize}
\item The generic finiteness statement producing a dense open of the target and an open of the source over which the restriction is finite, cited from the same chapter. Cited at proof lines 13345-13346. Provides the pair of opens in the separated case, which then has to be proved to exhaust the preimage.
\item A finite morphism is closed and proper, and the image of a proper scheme over the target is closed. Cited at proof lines 13347-13350. Shows that the open piece is also closed inside the preimage of the chosen open, which forces the complementary open to be empty.
\item A quasi-compact scheme contains a dense open separated subscheme, and for a quasi-compact morphism whose image of a closed subset misses a point there is an open neighbourhood of that point whose preimage misses that subset, both cited from the properties chapter and the same chapter. Cited at proof lines 13356-13364. Reduces the quasi-compact case to the separated case by first restricting to a separated dense open of the source.
\end{itemize}\medskip
\ifdefined\hypertarget\hypertarget{sourcepointer-1}{}\fi
\noindent\textit{Source reference.} Original lemma 01WN: focused-seeds/AGP/sources/stacks-snapshot-20261011/tags.txt, source lines 2488--2488. Same-chapter reference used in the retained author text at source line(s) 11282, 12507, 13352. It is supplied only as a navigation pointer so that the retained text compiles; the counted item delivers the author's own statement and the author's complete proof, and no part of the referenced item is copied in or substituted for the counted argument.\par\medskip
\ifdefined\hypertarget\hypertarget{sourcepointer-2}{}\fi
\noindent\textit{Source reference.} Original lemma 02NW: focused-seeds/AGP/sources/stacks-snapshot-20261011/tags.txt, source lines 3495--3495. Same-chapter reference used in the retained author text at source line(s) 13285, 13296, 13349, 13488, 13520. It is supplied only as a navigation pointer so that the retained text compiles; the counted item delivers the author's own statement and the author's complete proof, and no part of the referenced item is copied in or substituted for the counted argument.\par\medskip
\ifdefined\hypertarget\hypertarget{sourcepointer-3}{}\fi
\noindent\textit{Source reference.} Original lemma 01W6: focused-seeds/AGP/sources/stacks-snapshot-20261011/tags.txt, source lines 2471--2471. Same-chapter reference used in the retained author text at source line(s) 10104, 10627, 10794, 10905, 11283, 13351, 14146. It is supplied only as a navigation pointer so that the retained text compiles; the counted item delivers the author's own statement and the author's complete proof, and no part of the referenced item is copied in or substituted for the counted argument.\par\medskip
\ifdefined\hypertarget\hypertarget{sourcepointer-4}{}\fi
\noindent\textit{Source reference.} Original lemma 01RL: focused-seeds/AGP/sources/stacks-snapshot-20261011/tags.txt, source lines 2306--2306. Same-chapter reference used in the retained author text at source line(s) 1323, 1327, 1371, 13364. It is supplied only as a navigation pointer so that the retained text compiles; the counted item delivers the author's own statement and the author's complete proof, and no part of the referenced item is copied in or substituted for the counted argument.\par\medskip
\ifdefined\hypertarget\hypertarget{sourcepointer-5}{}\fi
\noindent\textit{Source reference.} Original lemma 03J1: focused-seeds/AGP/sources/stacks-snapshot-20261011/tags.txt, source lines 4623--4623. Same-chapter reference used in the retained author text at source line(s) 13360. It is supplied only as a navigation pointer so that the retained text compiles; the counted item delivers the author's own statement and the author's complete proof, and no part of the referenced item is copied in or substituted for the counted argument.\par\medskip
\begin{lemma}
\label{lemma-finite-over-dense-open}
Let $X$, $Y$ be schemes. Let $f : X \to Y$ be locally of finite type.
Let $X^0$, resp.\ $Y^0$ denote the set of generic points of irreducible
components of $X$, resp.\ $Y$. Assume
\begin{enumerate}
\item $X^0$ and $Y^0$ are finite and $f^{-1}(Y^0) = X^0$,
\item either $f$ is quasi-compact or $f$ is separated.
\end{enumerate}
Then there exists a dense open $V \subset Y$
such that $f^{-1}(V) \to V$ is finite.
\end{lemma}

\begin{proof}
Since $Y$ has finitely many irreducible components, we can find a dense
open which is a disjoint union of its irreducible components. Thus we may
assume $Y$ is irreducible affine with generic point $\eta$. Then the fibre
over $\eta$ is finite as $X^0$ is finite.

\medskip\noindent
Assume $f$ is separated and $Y$ irreducible affine. Choose $V \subset Y$
and $U \subset X$ as in Lemma \ref{lemma-generically-finite} part (3).
Since $f|_U : U \to V$ is finite, we see that $U \subset f^{-1}(V)$
is closed as well as open (Lemmas \ref{lemma-image-proper-scheme-closed} and
\ref{lemma-finite-proper}). Thus $f^{-1}(V) = U \amalg W$ for some
open subscheme $W$ of $X$. However, since $U$ contains all the generic
points of $X$ we conclude that $W = \emptyset$ as desired.

\medskip\noindent
Assume $f$ is quasi-compact and $Y$ irreducible affine. Then $X$ is
quasi-compact, hence there exists a dense open subscheme $U \subset X$
which is separated
(Properties, Lemma \ref{properties-lemma-quasi-compact-dense-open-separated}).
Since the set of generic points $X^0$ is finite, we see that $X^0 \subset U$.
Thus $\eta \not \in f(X \setminus U)$. Since $X \setminus U \to Y$ is
quasi-compact, we conclude that there is a nonempty open $V \subset Y$
such that $f^{-1}(V) \subset U$, see Lemma \ref{lemma-quasi-compact-dominant}.
After replacing $X$ by $f^{-1}(V)$ and $Y$ by $V$ we reduce to the
separated case which we dealt with in the preceding paragraph.
\end{proof}

\end{document}
